\documentclass[a4paper, serif, 10pt]{moderncv}

%% moderncv
% style and color
\moderncvstyle{banking}
\moderncvcolor{black}

% renew moderncv command to adapt for CJK colon
\renewcommand*{\cvitem}[3][.25em]{%
  \ifstrempty{#2}{}{\hintstyle{#2}：}{#3}%
  \par\addvspace{#1}}

\renewcommand*{\cvitemwithcomment}[4][.25em]{%
  \savebox{\cvitemwithcommentmainbox}{\ifstrempty{#2}{}{\hintstyle{#2}：}#3}%
  \setlength{\cvitemwithcommentmainlength}{\widthof{\usebox{\cvitemwithcommentmainbox}}}%
  \setlength{\cvitemwithcommentcommentlength}{\maincolumnwidth-\separatorcolumnwidth-\cvitemwithcommentmainlength}%
  \begin{minipage}[t]{\cvitemwithcommentmainlength}\usebox{\cvitemwithcommentmainbox}\end{minipage}%
  \hfill% fill of \separatorcolumnwidth
  \begin{minipage}[t]{\cvitemwithcommentcommentlength}\raggedleft\small\itshape#4\end{minipage}%
  \par\addvspace{#1}}

%% page layout/margins
\usepackage[top=2.5cm, bottom=2.5cm, left=1.5cm, right=1.5cm]{geometry}

\nopagenumbers{}

%% Babel config for English language
\usepackage[english]{babel}

%% fontspec
\usepackage{fontspec}

\IfFontExistsTF{Linux Libertine}{
  \setmainfont[Ligatures={TeX, Common}, Numbers=Lining, ItalicFont=Linux Libertine]{Linux Libertine}
}{}
\IfFontExistsTF{Linux Libertine O}{
  \setmainfont[Ligatures={TeX, Common}, Numbers=Lining, ItalicFont=Linux Libertine O]{Linux Libertine O}
}{}

%% CTeX
% CJK support, used to show CJK characters in the resume
%
% - fontset=none: disable builtin fontset but instead set the CJK font manually
% - heading=false: disable ctex heading
% - punct=kaiming: use kaiming punctuations styles for CJK
% - scheme=plain: use plain scheme, do not override \normalsize font size
% - space=auto: space settings for CJK characters
%
% ref:
% - http://ctan.mirrorcatalogs.com/language/chinese/ctex/ctex.pdf
\usepackage[UTF8, heading=false, punct=kaiming, scheme=plain, space=auto]{ctex}

\IfFontExistsTF{Noto Serif CJK SC}{
  \setCJKmainfont{Noto Serif CJK SC}
}{}
\IfFontExistsTF{Noto Sans CJK SC}{
  \setCJKsansfont{Noto Sans CJK SC}
}{}

%% line spacing
\usepackage{setspace}
\setstretch{1.125}

%% URL styling - use normal text instead of monospace
\urlstyle{same}

% Auto-underline all links
% Defer until \begin{document} so hyperref is already loaded
\AtBeginDocument{%
  \let\oldhref\href
  \renewcommand{\href}[2]{\oldhref{#1}{\underline{#2}}}%
}

%% Patch \httplink and \httpslink to handle full URLs
% The original moderncv macros always prepend http:// or https:// to URLs.
% This causes issues when the URL already has a protocol (e.g., https://example.com)
% resulting in malformed URLs like https://https://example.com.
% This patch checks if the URL already contains "://" and uses it directly if so.
% Note: We use \str_set:Nx (x-expansion) to expand macros like \@homepage first.
\ExplSyntaxOn
\str_new:N \g_yamlresume_url_str

\RenewDocumentCommand{\httpslink}{O{}m}{
  \str_set:Nx \g_yamlresume_url_str {#2}
  \str_if_in:NnTF \g_yamlresume_url_str {://}
    { \url{#2} }
    { \url{https://#2} }
}

\RenewDocumentCommand{\httplink}{O{}m}{
  \str_set:Nx \g_yamlresume_url_str {#2}
  \str_if_in:NnTF \g_yamlresume_url_str {://}
    { \url{#2} }
    { \url{http://#2} }
}
\ExplSyntaxOff

\name{佐藤 美咲}{}
\title{カスタマーサクセスマネージャー | SaaS導入支援と顧客成長の推進}
\phone[mobile]{+81 90-1234-5678}
\email{misaki.sato@example.com}
\homepage{https://example.com/misaki-sato}

\address{日本、東京都、渋谷区、東京都渋谷区神南1-2-3、150-0001}{}{}

\extrainfo{{\small \faLinkedin}\ \href{https://www.linkedin.com/in/misaki-sato-csm}{@misaki-sato-csm} {} {} {} • {} {} {} 
{\small \faTwitter}\ \href{https://twitter.com/misaki\_csm}{@misaki\_csm}}

\begin{document}

\maketitle

\section{基本情報}

\cvline{}{\begin{itemize}
\item SaaS企業で5年以上のカスタマーサクセス経験を持ち、導入支援から定着支援、更新・アップセルまで一貫して担当
\item 顧客のビジネス目標を理解し、プロダクト活用を通じた成果創出を支援
\item データ分析に基づくヘルススコア設計とチャーン防止施策の立案・実行が強み
\item クロスファンクショナルなチームと連携し、顧客の声をプロダクト改善に還元
\item 日本語と英語でのコミュニケーションが可能で、グローバル顧客にも対応
\end{itemize}}
\section{学歴}

\cventry{2014年4月～2018年3月}
        {学士、経営学、成績：3.7/4.0}
        {早稲田大学}
        {\url{https://www.waseda.jp}}
        {}
        {\begin{itemize}
\item 経営学を専攻し、マーケティングと組織マネジメントを中心に学習
\item データ分析や消費者行動の授業を通じて、顧客理解の基礎を習得
\item グループワークやプレゼンテーションで協働と発信力を養成
\end{itemize}
\textbf{コース}：マーケティング戦略、経営組織論、消費者行動論、データ分析入門、ビジネスコミュニケーション、国際経営}
\section{職歴}

\cventry{2021年4月～現在}
        {シニアカスタマーサクセスマネージャー}
        {株式会社クラウドコネクト}
        {\url{https://cloudconnect.example.jp}}
        {}
        {\begin{itemize}
\item エンタープライズ顧客30社以上を担当し、導入支援から定着、更新、アップセルまで一貫して推進
\item 顧客のKPIに基づく成功計画を策定し、四半期ビジネスレビュー（QBR）を主導
\item ヘルススコアを再設計し、チャーンリスクの早期検知と proactive な支援を実現
\item プロダクトチームと連携し、顧客要望を機能改善に反映
\item 年間契約更新率を92\%から96\%へ改善し、アップセル収益を前年比130\%に拡大
\end{itemize}
\textbf{キーワード}：カスタマーサクセス、導入支援、更新管理、アップセル、QBR、チャーン防止}

\cventry{2018年4月～2021年3月}
        {カスタマーサクセスマネージャー}
        {株式会社サポートテック}
        {\url{https://supporttech.example.jp}}
        {}
        {\begin{itemize}
\item 中小規模顧客100社以上を担当し、オンボーディングと活用支援を実施
\item 顧客向けトレーニングやウェビナーを企画・運営し、プロダクト利用率を向上
\item 顧客満足度調査を分析し、サポート品質改善とナレッジベース整備に貢献
\item 解約予兆のある顧客へのリテンション施策を立案し、更新率を88\%から93\%へ改善
\end{itemize}
\textbf{キーワード}：オンボーディング、活用支援、リテンション、トレーニング、顧客満足度}

\cventry{2016年4月～2018年3月}
        {カスタマーサポート担当}
        {株式会社グローバルソリューションズ}
        {\url{https://globalsolutions.example.jp}}
        {}
        {\begin{itemize}
\item 法人顧客からの問い合わせ対応を通じて、プロダクトの機能理解と課題解決を支援
\item 問い合わせデータを分析し、FAQやマニュアルの改善を提案
\item サポートチーム内のナレッジ共有を促進し、対応時間の短縮に貢献
\end{itemize}
\textbf{キーワード}：カスタマーサポート、問い合わせ対応、ナレッジ管理、課題解決}
\section{言語}

\cvline{日本語}{ネイティブ・母国語レベル \hfill \textbf{キーワード}：母国語、ビジネス文書、プレゼンテーション}
\cvline{英語}{完全な実務レベル \hfill \textbf{キーワード}：TOEIC 920、ビジネス会話、英文メール、海外顧客対応}
\section{スキル}

\cvline{カスタマーサクセス}{エキスパート \hfill \textbf{キーワード}：オンボーディング、導入支援、顧客関係構築、チャーン防止、アップセル、更新管理}
\cvline{データ分析}{上級 \hfill \textbf{キーワード}：Salesforce、Tableau、SQL、Excel、KPI設計、ヘルススコア}
\cvline{プロジェクト管理}{上級 \hfill \textbf{キーワード}：Asana、Jira、Notion、Slack、進行管理}
\cvline{コミュニケーション}{エキスパート \hfill \textbf{キーワード}：プレゼンテーション、折衝、ファシリテーション、英語、日本語}
\section{受賞・表彰}

\cventry{2023年12月}
        {株式会社クラウドコネクト}
        {年間最優秀カスタマーサクセスマネージャー}
        {}
        {}
        {顧客満足度と更新率の向上、およびアップセル収益の拡大に貢献した功績により表彰。}

\cventry{2020年12月}
        {株式会社サポートテック}
        {社長賞}
        {}
        {}
        {大手顧客のチャーン防止プロジェクトをリードし、年間契約更新率を93\%まで改善した功績により受賞。}
\section{資格・免状}

\cventry{2022年6月}
        {Salesforce}
        {Salesforce 認定アドミニストレーター}
        {\url{https://trailhead.salesforce.com/credentials/administrator}}
        {}
        {}

\cventry{2021年9月}
        {HubSpot}
        {HubSpot インバウンドセールス認定}
        {\url{https://academy.hubspot.com/certification}}
        {}
        {}

\cventry{2020年3月}
        {AXELOS}
        {ITIL 4 Foundation}
        {\url{https://www.axelos.com/certifications/itil-certifications}}
        {}
        {}
\section{出版・著作}

\cventry{2023年8月}
        {SaaSにおけるカスタマーサクセス指標の設計}
        {テックビジネスジャーナル}
        {\url{https://example.com/articles/customer-success-kpi}}
        {}
        {\begin{itemize}
\item カスタマーサクセスにおけるKPI設計の考え方とヘルススコア活用事例を解説
\item チャーン予測とアップセル機会の可視化による収益貢献の可能性を提示
\item 顧客セグメント別の成功指標の違いと運用上の注意点を整理
\end{itemize}}
\section{プロジェクト}

\cventry{2023年1月～2023年12月}
        {顧客データを活用したヘルススコアモデルの再設計}
        {顧客ヘルススコア改善プロジェクト}
        {\url{https://example.com/projects/health-score}}
        {}
        {\begin{itemize}
\item 利用頻度、サポート問い合わせ、NPSなどのデータを統合し、チャーン予測精度を向上
\item リスク顧客に対する自動アラートと担当者アクションの標準化を実施
\item プロジェクト完了後、解約率を前年比20\%削減
\end{itemize}
\textbf{キーワード}：ヘルススコア、データ分析、チャーン予測、顧客セグメント}

\cventry{2022年4月～2022年9月}
        {新規顧客向けオンボーディングプロセスの自動化}
        {オンボーディング自動化プロジェクト}
        {\url{https://example.com/projects/onboarding-automation}}
        {}
        {\begin{itemize}
\item 契約後から初回価値実感までのステップを整理し、自動化ツールを導入
\item テンプレート化により、担当者ごとのばらつきを削減
\item オンボーディング完了までの期間を平均14日から7日へ短縮
\end{itemize}
\textbf{キーワード}：オンボーディング、自動化、プロセス改善、顧客導入}
\section{興味・関心}

\cvline{読書}{ビジネス書、心理学、テクノロジー}
\cvline{旅行}{国内旅行、海外旅行、温泉}
\cvline{料理}{イタリア料理、和食、パン作り}
\section{ボランティア活動}

\cventry{2019年9月～現在}
        {メンター}
        {一般社団法人 キャリア支援ネットワーク}
        {\url{https://example.com/volunteer}}
        {}
        {\begin{itemize}
\item 若手社会人や学生向けに、キャリア相談とビジネススキル向上のメンタリングを実施
\item 月1回のワークショップを企画し、コミュニケーションや顧客対応の基礎を指導
\item メンティーの成長を支援し、延べ50名以上を担当
\end{itemize}}
    
\end{document}